\section{Introduction}\label{sec:introduction}

A public compatibility standard can benefit consumers and firms by placing products on a common technical interface. Yet agreement about the desirability of compatibility need not produce a common standard. Governments care about the profits of domestic producers as well as the benefits to consumers, and these interests may differ across countries. An additional institutional obstacle is easy to overlook: a proposal cannot pass a vote unless someone with the legal right to place it on the agenda chooses to do so. A standard may have enough supporters to win, but no authorized sponsor willing to introduce it.

The agenda question becomes more consequential when firms choose their technologies before the government vote. A firm investing in a partially compatible architecture cannot later recover that expenditure simply because a public interface becomes mandatory. The expectation of future intervention therefore changes what firms are willing to develop in the first place. Conversely, an observed private regional standard may give an outsider government a reason to block the public alternative even though the two founding countries favor it. To evaluate the resulting equilibrium, one must connect the initial private investment decisions to credible national decisions after \emph{every} possible technological history.

We develop a three-country oligopoly model that isolates this interaction. Each country hosts one firm and one government. Firms compete in differentiated-product Cournot markets and simultaneously decide whether to develop an open technical interface. A developer selects two continuous attributes: compatibility capacity and commonality of design. Effective private compatibility is limited by what both firms have developed, so a bilateral standard depends on the minimum investment by its two founders. Technical choices increase demand or change competitive substitution, but entail sunk setup and convex physical costs. The design space is continuous and firms choose independently; technologies need not be symmetric.

Once investments are observed, governments decide whether to mandate a single legally available public interface. The constitution specifies both which national governments can \emph{propose} a mandate and how many affirmative votes it needs. Under a successful mandate, firms decide whether to comply with the public interface or cease production. A developer in the same open technological family pays only the cost of completing its missing compatibility stock. Its original commonality investment, however, is sunk and cannot be recovered. We solve the ensuing compliance, voting, and initial investment games in pure strategies. The public mandate is enforceable, but this is not a model of compulsory licensing, patent expropriation, or governments bargaining over the design of an otherwise unavailable technology.

Our first result identifies a sharp constitutional threshold in an explicit region of economic primitives. When one or two affirmative national votes are sufficient, \emph{public standardization is the only pure-strategy subgame-perfect equilibrium terminal outcome if and only if at least two governments hold independent proposal rights}. With a sole proposal holder, the other two countries can sustain a privately developed bilateral standard, and the public standard is also supported by an equilibrium. Unanimity preserves privately developed regional equilibria regardless of how proposal rights are distributed. The threshold thus concerns the set of \emph{equilibrium outcomes}, not the uniqueness of each initial technological choice. The analysis does not rule out mixed-strategy equilibria.

Why does this institutional threshold extend beyond voting over an assumed private standard? In a three-country polity, any two authorized proposers intersect every two-government majority. The observation is elementary. But an arbitrary private technology profile need not yield two favorable votes. Some heterogeneous private standards can indeed make particular governments oppose a public interface. Our central equilibrium task is therefore to rule out private investment profiles that would both evade the public vote \emph{and} survive a firm's incentive to deviate. With two developers, excessive unilateral compatibility or commonality capital can be reduced without changing the common pair technology; at higher commonality choices, a smaller design adjustment raises private profit while leaving political nonintervention intact. With three developers, surviving private profiles must have ties in their maximal design attributes; careful analysis of the remaining asymmetric cases shows that two national governments weakly support the public mandate. Only after these private-investment restrictions are established does the institutional majority-intersection argument complete the result. Our equilibrium claim is consequently stronger than an ex-post vote calculation at a symmetric regional benchmark.

The second result concerns welfare and the timing of the mandate. We distinguish three thresholds evaluated at the privately optimal two-firm regional design. The excluded government's \emph{veto threshold} measures when it individually supports mandatory public standardization. The \emph{ex-post welfare threshold} asks whether society gains by requiring a retrofit \emph{after} the regional firms have already made their investments. The \emph{ex-ante welfare threshold} instead compares building and keeping the regional design with selecting the public interface from the start. We derive a strict ordering of these thresholds. In particular, the last two differ by exactly the cost of the private commonality investment that becomes redundant under the public interface.

That wedge has a direct interpretation. A public interface can dominate private regional development when chosen \emph{before} investment, yet it can be costly to impose after firms have sunk resources into the alternative. There is also a region in which the privately developed regional equilibrium yields higher ex-ante world welfare than the public equilibrium that becomes uniquely available when proposal rights are broadened. Thus a constitutional reform that eliminates private regional equilibria does not necessarily increase aggregate welfare. When the initial constitution supports both public and regional equilibria, an overall welfare comparison still requires an equilibrium-selection rule; we do not infer a unique realized gain from the change in the equilibrium set.

Several established lines of research bear directly on this question. \citet{KatzShapiro1985} and \citet{FarrellSaloner1988} show why private compatibility incentives and coordination institutions matter. \citet{GandalShy2001} already examines foreign-standard recognition and the formation of potentially exclusionary three-country standardization unions. \citet{Klimenko2009} shows how international governments' strategic compatibility policies can produce too much or too little compatibility, while \citet{Costinot2008} compares national treatment and mutual recognition and obtains welfare rankings that depend on the externality environment. We therefore cannot present international standardization coalitions, governmental conflicts, or a reversal of welfare rankings as new in themselves.

The closest institutional and investment ideas likewise have well-established antecedents. \citet{FarrellSimcoe2012} analyze consensus standardization with neutral intervention and endogenous ex-ante improvement incentives; their discussion explicitly considers a two-vote rule in a different governance game. \citet{Ruiz2004} examines government recognition policy alongside endogenous product characteristics. Our narrower question is whether independently held \emph{national proposal rights}, combined with a voting quorum, eliminate \emph{all privately sustainable pure-SPE terminal regional standards} when each firm can choose continuous \emph{sunk} open-technology stocks. The difference lies in the sequential equilibrium domain and the regulatory instrument, not in a claim to invent voting or product differentiation in standardization research.

The model is intentionally restrictive. Public \(O(1,0)\) is preexisting and legally accessible, private interfaces belong to the open family, and governments cannot make compensating monetary transfers or select a new public standard. The main result holds on a stated rectangle of the primitives and an additional conservative local robustness neighborhood. Those assumptions make the contribution testable but limit its interpretation. The full published proposition and proof comparisons with some of the closest earlier international-policy work also remain part of the final pre-submission novelty audit; we make no claim that the complete literature has been ruled out as a source of a more general covering theorem.

The paper proceeds as follows. Section~\ref{sec:model} defines the private technology stocks, demand system, and political sequence. Section~\ref{sec:policy} solves the public-compliance subgame and derives the government policy gains. Section~\ref{sec:equilibria} characterizes the institutional threshold over the entire continuous investment domain. Section~\ref{sec:welfare} separates national veto, ex-post retrofit efficiency, and ex-ante allocation efficiency. Section~\ref{sec:robustness} discusses the jointly perturbed primitive neighborhood, earlier literature, and the assumptions carrying the result. The appendices make the proof-critical derivations and rational certificate regions explicit.
